\documentclass[11pt,a4paper]{article}
\usepackage[T1]{fontenc}
\usepackage[utf8]{inputenc}
\usepackage[french]{babel}
\usepackage{geometry,xcolor,listings,hyperref,amsmath}
\geometry{margin=23mm}
\definecolor{navy}{RGB}{18,49,88}
\lstset{language=Python,basicstyle=\ttfamily\small,breaklines=true,breakatwhitespace=false,
  frame=single,backgroundcolor=\color{black!3},columns=fullflexible,
  showstringspaces=false,keepspaces=true}
\hypersetup{colorlinks=true,linkcolor=navy}
\title{MLP et fonction sinus : explication ligne par ligne}
\author{Frédéric Magoulès}
\date{November 2026}
\begin{document}
\maketitle
\noindent Ce document commente, sans modifier le programme, chaque ligne non vide du fichier \texttt{mlp\_sine\_intro.py}. Les numéros correspondent exactement au fichier fourni. Les lignes vides ne sont pas commentées. La validation est réalisée après chaque époque ; l'inférence finale réutilise les abscisses de validation, sans jeu de test séparé.
\section*{Explications détaillées}
\subsection*{Ligne 1}
\begin{lstlisting}
"""Introductory MLP: fit sin(2*pi*x) with PyTorch on CPU."""
\end{lstlisting}
Docstring du module : annonce l'objectif, l'approximation de sin(2$\pi$x), PyTorch et l'exécution sur CPU.

\subsection*{Ligne 2}
\begin{lstlisting}
from pathlib import Path
\end{lstlisting}
Importe Path pour construire un chemin de sortie indépendant du système d'exploitation.

\subsection*{Ligne 3}
\begin{lstlisting}
import numpy as np
\end{lstlisting}
Importe NumPy, utilisé pour fixer la graine aléatoire.

\subsection*{Ligne 4}
\begin{lstlisting}
import torch
\end{lstlisting}
Importe PyTorch, bibliothèque de calcul tensoriel et de différentiation automatique.

\subsection*{Ligne 5}
\begin{lstlisting}
from torch import nn
\end{lstlisting}
Importe les couches et fonctions de réseaux de neurones de PyTorch sous le nom nn.

\subsection*{Ligne 6}
\begin{lstlisting}
from torch.utils.data import DataLoader, TensorDataset
\end{lstlisting}
Importe les outils de constitution d'un jeu de données et de mini-lots.

\subsection*{Ligne 7}
\begin{lstlisting}
import matplotlib
\end{lstlisting}
Importe Matplotlib, bibliothèque de visualisation.

\subsection*{Ligne 8}
\begin{lstlisting}
matplotlib.use("Agg")
\end{lstlisting}
Choisit le moteur de rendu non interactif Agg, avant d'importer pyplot ; cela permet l'enregistrement de figures sans interface graphique.

\subsection*{Ligne 9}
\begin{lstlisting}
import matplotlib.pyplot as plt
\end{lstlisting}
Importe l'interface de tracé pyplot.

\subsection*{Ligne 11}
\begin{lstlisting}
torch.manual_seed(7)
\end{lstlisting}
Fixe la graine pseudo-aléatoire de PyTorch à 7 pour rendre l'initialisation et l'ordre aléatoire plus reproductibles.

\subsection*{Ligne 12}
\begin{lstlisting}
np.random.seed(7)
\end{lstlisting}
Fixe également la graine pseudo-aléatoire de NumPy.

\subsection*{Ligne 14}
\begin{lstlisting}
# Training and validation inputs are distinct.
\end{lstlisting}
Commentaire : les abscisses d'apprentissage et de validation sont distinctes.

\subsection*{Ligne 15}
\begin{lstlisting}
x_train = torch.linspace(0, 1, 101).unsqueeze(1)
\end{lstlisting}
Construit 101 abscisses uniformément réparties entre 0 et 1 ; unsqueeze(1) transforme le vecteur en tenseur colonne de forme (101,1).

\subsection*{Ligne 16}
\begin{lstlisting}
x_val = torch.linspace(0.005, 0.995, 100).unsqueeze(1)
\end{lstlisting}
Construit 100 abscisses de validation entre 0,005 et 0,995 ; elles sont décalées par rapport à celles d'apprentissage.

\subsection*{Ligne 17}
\begin{lstlisting}
y_train = torch.sin(2 * torch.pi * x_train)
\end{lstlisting}
Calcule les valeurs cibles sin(2$\pi$x) pour les points d'apprentissage.

\subsection*{Ligne 18}
\begin{lstlisting}
y_val = torch.sin(2 * torch.pi * x_val)
\end{lstlisting}
Calcule les mêmes valeurs exactes pour les points de validation.

\subsection*{Ligne 19}
\begin{lstlisting}
loader = DataLoader(TensorDataset(x_train, y_train), batch_size=32, shuffle=True)
\end{lstlisting}
Regroupe abscisses et cibles, puis crée des mini-lots de taille maximale 32, mélangés à chaque parcours.

\subsection*{Ligne 21}
\begin{lstlisting}
# One input, one output, and two small hidden layers.
\end{lstlisting}
Commentaire : le réseau reçoit un scalaire et prédit un scalaire, avec deux couches cachées.

\subsection*{Ligne 22}
\begin{lstlisting}
model = nn.Sequential(nn.Linear(1, 32), nn.Tanh(), nn.Linear(32, 32),
\end{lstlisting}
Commence un réseau séquentiel : première couche affine 1$\to$32, puis activation tanh, puis deuxième couche affine 32$\to$32.

\subsection*{Ligne 23}
\begin{lstlisting}
                      nn.Tanh(), nn.Linear(32, 1))
\end{lstlisting}
Ajoute la seconde activation tanh et la couche de sortie affine 32$\to$1 ; termine la définition du réseau.

\subsection*{Ligne 24}
\begin{lstlisting}
optimizer = torch.optim.Adam(model.parameters(), lr=1e-2)
\end{lstlisting}
Crée l'optimiseur Adam sur les paramètres du réseau, avec un taux d'apprentissage de 0,01.

\subsection*{Ligne 25}
\begin{lstlisting}
loss_fn = nn.MSELoss()
\end{lstlisting}
Choisit la fonction de perte MSE : moyenne des carrés des écarts entre prédictions et cibles.

\subsection*{Ligne 26}
\begin{lstlisting}
train_loss, val_loss = [], []
\end{lstlisting}
Initialise deux listes pour enregistrer, à chaque époque, les pertes d'apprentissage et de validation.

\subsection*{Ligne 28}
\begin{lstlisting}
# Training
\end{lstlisting}
Marque le début de la phase d'apprentissage.

\subsection*{Ligne 29}
\begin{lstlisting}
for epoch in range(200):
\end{lstlisting}
Répète l'apprentissage pendant 200 époques (un passage complet sur les données par époque).

\subsection*{Ligne 30}
\begin{lstlisting}
    model.train()
\end{lstlisting}
Place le modèle en mode apprentissage ; utile notamment pour dropout et batch normalization, absents ici.

\subsection*{Ligne 31}
\begin{lstlisting}
    for xb, yb in loader:
\end{lstlisting}
Parcourt les mini-lots fournis par DataLoader ; xb contient les entrées et yb leurs cibles.

\subsection*{Ligne 32}
\begin{lstlisting}
        prediction = model(xb)           # Forward: predict y from x.
\end{lstlisting}
Passe xb dans le MLP pour calculer la prédiction : propagation avant (forward).

\subsection*{Ligne 33}
\begin{lstlisting}
        loss = loss_fn(prediction, yb)   # Loss: compare to sin(2*pi*x).
\end{lstlisting}
Calcule l'erreur quadratique moyenne sur le mini-lot courant.

\subsection*{Ligne 34}
\begin{lstlisting}
        optimizer.zero_grad()           # Clear gradients from previous batch.
\end{lstlisting}
Remet à zéro les gradients précédemment accumulés par l'optimiseur.

\subsection*{Ligne 35}
\begin{lstlisting}
        loss.backward()                 # Backward: compute gradients.
\end{lstlisting}
Effectue la rétropropagation : calcule les dérivées de la perte par rapport aux paramètres.

\subsection*{Ligne 36}
\begin{lstlisting}
        optimizer.step()                # Update the MLP parameters.
\end{lstlisting}
Effectue une mise à jour des poids et biais à l'aide d'Adam.

\subsection*{Ligne 37}
\begin{lstlisting}
    # Validation
\end{lstlisting}
Marque le début de la validation effectuée après chaque époque.

\subsection*{Ligne 38}
\begin{lstlisting}
    model.eval()
\end{lstlisting}
Place le modèle en mode évaluation.

\subsection*{Ligne 39}
\begin{lstlisting}
    with torch.no_grad():
\end{lstlisting}
Désactive l'enregistrement des opérations pour le calcul des gradients : évaluation moins coûteuse.

\subsection*{Ligne 40}
\begin{lstlisting}
        train_loss.append(loss_fn(model(x_train), y_train).item())
\end{lstlisting}
Calcule la MSE sur tous les points d'apprentissage et stocke sa valeur scalaire Python.

\subsection*{Ligne 41}
\begin{lstlisting}
        val_loss.append(loss_fn(model(x_val), y_val).item())
\end{lstlisting}
Calcule la MSE sur les points distincts de validation et la stocke.

\subsection*{Ligne 42}
\begin{lstlisting}
    if (epoch + 1) % 40 == 0:
\end{lstlisting}
Teste si le numéro d'époque (compté à partir de 1) est un multiple de 40.

\subsection*{Ligne 43}
\begin{lstlisting}
        print(f"Epoch {epoch+1:3d}: train={train_loss[-1]:.6f}, val={val_loss[-1]:.6f}")
\end{lstlisting}
Affiche périodiquement les pertes d'apprentissage et de validation avec six décimales.

\subsection*{Ligne 45}
\begin{lstlisting}
# Inference
\end{lstlisting}
Marque la phase d'inférence finale : les paramètres appris ne sont plus mis à jour.

\subsection*{Ligne 46}
\begin{lstlisting}
model.eval()
\end{lstlisting}
Place le modèle en mode évaluation.

\subsection*{Ligne 47}
\begin{lstlisting}
with torch.no_grad():
\end{lstlisting}
Désactive le suivi des gradients pendant l'inférence.

\subsection*{Ligne 48}
\begin{lstlisting}
    prediction = model(x_val).numpy().ravel()
\end{lstlisting}
Prédit les valeurs sur x\_val, convertit le tenseur en tableau NumPy et l'aplatit en un vecteur 1D.

\subsection*{Ligne 50}
\begin{lstlisting}
fig, axes = plt.subplots(1, 2, figsize=(10, 3.8))
\end{lstlisting}
Crée une figure avec deux sous-graphiques côte à côte, de taille 10 $\times$ 3,8 pouces.

\subsection*{Ligne 51}
\begin{lstlisting}
axes[0].semilogy(range(1, 201), train_loss, label="Training")
\end{lstlisting}
Trace l'historique de la perte d'apprentissage en ordonnée logarithmique.

\subsection*{Ligne 52}
\begin{lstlisting}
axes[0].semilogy(range(1, 201), val_loss, label="Validation")
\end{lstlisting}
Trace l'historique de la perte de validation sur les mêmes axes.

\subsection*{Ligne 53}
\begin{lstlisting}
axes[0].set(xlabel="Epoch", ylabel="Mean squared error", title="Training history")
\end{lstlisting}
Définit les noms des axes et le titre du premier graphique.

\subsection*{Ligne 54}
\begin{lstlisting}
axes[0].legend()
\end{lstlisting}
Affiche la légende du premier graphique.

\subsection*{Ligne 55}
\begin{lstlisting}
axes[1].plot(x_val.numpy().ravel(), y_val.numpy().ravel(), label="Exact sine")
\end{lstlisting}
Trace la fonction sinus exacte sur les abscisses de validation.

\subsection*{Ligne 56}
\begin{lstlisting}
axes[1].plot(x_val.numpy().ravel(), prediction, "--", label="MLP prediction")
\end{lstlisting}
Trace les prédictions du MLP en pointillés, sur les mêmes abscisses.

\subsection*{Ligne 57}
\begin{lstlisting}
axes[1].set(xlabel="x", ylabel="y", title="Function approximation")
\end{lstlisting}
Définit les noms des axes et le titre du deuxième graphique.

\subsection*{Ligne 58}
\begin{lstlisting}
axes[1].legend()
\end{lstlisting}
Affiche la légende du deuxième graphique.

\subsection*{Ligne 59}
\begin{lstlisting}
fig.tight_layout()
\end{lstlisting}
Ajuste automatiquement les espacements pour limiter les chevauchements.

\subsection*{Ligne 60}
\begin{lstlisting}
output = Path(__file__).with_name("mlp_sine_results.png") if "__file__" in globals() else Path("mlp_sine_results.png")
\end{lstlisting}
Détermine le chemin de sortie : à côté du script si \_\_file\_\_ existe, sinon dans le répertoire courant (cas fréquent en notebook).

\subsection*{Ligne 61}
\begin{lstlisting}
fig.savefig(output, dpi=160)
\end{lstlisting}
Enregistre la figure PNG avec une résolution de 160 points par pouce.

\subsection*{Ligne 62}
\begin{lstlisting}
print(f"Saved {output}")
\end{lstlisting}
Affiche le chemin du fichier image enregistré.

\subsection*{Ligne 63}
\begin{lstlisting}
plt.show()
\end{lstlisting}
Affiche la figure lorsque l'environnement graphique le permet.

\end{document}
